% ----------------------------------------------------------------------
%  2.2.6. Дислокации в типичных кристаллических структурах
% ----------------------------------------------------------------------
\subsection{Дислокации в типичных кристаллических структурах}\label{ssec:disl-in-structures}

\begin{orientir}
До сих пор мы рисовали дислокации в простой кубической решётке. Реальные
металлы имеют решётки ГЦК, ОЦК и ГПУ, и в них появляются новые объекты:
\textbf{частичные дислокации}, после прохождения которых решётка не
восстанавливается полностью, а остаётся \textbf{дефект упаковки}. Здесь
мы разберёмся, какие дислокации энергетически выгодны (критерий
Франка), как устроены плотнейшие упаковки и дефекты в них, и что такое
частичные дислокации Шокли и Франка.
\end{orientir}

\subsubsection{Полные и частичные дислокации}

\paragraph{Простая кубическая решётка}
Пробег дислокации с вектором Бюргерса, равным вектору решётки,
восстанавливает расположение атомов в зоне сдвига: решётка совмещается
сама с собой (\term{тождественная трансляция}). В простой кубической
решётке тождественную трансляцию дают векторы $a\nfam{100}$ (ребро куба),
$a\nfam{110}$ (диагональ грани) и $a\nfam{111}$ (диагональ куба).
Дислокации с такими векторами называют \term{единичными}. Возможны и
дислокации \emph{$n$-кратной мощности} с вектором $n\vb b$.

Но энергия дислокации пропорциональна $b^2$. Поэтому:
\begin{itemize}
\item дислокация $n$-кратной мощности неустойчива и распадается на $n$
единичных: $n^2b^2>n\cdot b^2$. Например, если у краевой дислокации две
экстраплоскости оканчиваются на одном уровне, они «разбегаются», и
получаются две отдельные дислокации;
\item из единичных дислокаций преобладают дислокации с \emph{наименьшим}
вектором Бюргерса (в простой кубической~— $a\nfam{100}$).
\end{itemize}

\begin{opredelenie}
\term{Полная дислокация}~— дислокация, вектор Бюргерса которой равен
вектору трансляции решётки; её пробег оставляет за собой совершенную
решётку. \term{Частичная} (неполная) \term{дислокация}~— дислокация,
вектор Бюргерса которой \emph{не} является вектором трансляции; после её
пробега атомы оказываются в механически устойчивых, но «неправильных»
положениях~— образуется дефект упаковки.
\end{opredelenie}

В ГЦК, ОЦК и ГПУ структурах вектор Бюргерса частичной дислокации обычно
меньше, чем у полной, и её энергия тоже меньше. Частичные дислокации
играют исключительно важную роль во многих процессах.

\begin{vnimanie}
Деление на краевые, винтовые и смешанные и деление на полные и частичные
основаны на \emph{разных} признаках. Первое~— на угле между линией и
$\vb b$, второе~— на величине $\vb b$ по сравнению с вектором трансляции.
Полная дислокация может быть краевой, винтовой или смешанной; смешанная~—
полной или частичной.
\end{vnimanie}

\subsubsection{Энергетический критерий дислокационных реакций}

Дислокации могут вступать в реакции: полная может расщепиться на
частичные ($\vb b_1=\vb b_2+\vb b_3$), частичные~— объединиться
($\vb b_1+\vb b_2=\vb b_3$) или перекомбинироваться
($\vb b_1+\vb b_2=\vb b_3+\vb b_4$). Слева записывают векторы Бюргерса
исходных дислокаций, справа~— продуктов.

\paragraph{Закон сохранения}
Сумма векторов Бюргерса до реакции равна сумме после. Для реакции
$k_1[u_1v_1w_1]=k_2[u_2v_2w_2]+k_3[u_3v_3w_3]$ это значит, что совпадают
все три компоненты: $k_1u_1=k_2u_2+k_3u_3$, $k_1v_1=k_2v_2+k_3v_3$,
$k_1w_1=k_2w_2+k_3w_3$.

\paragraph{Критерий Франка}
Из всех реакций, разрешённых законом сохранения, идут только те, что
уменьшают энергию:
\begin{keyeq}\label{eq:frank-crit}
\text{реакция идёт, если}\qquad \sum_{\text{до}}b_i^2>\sum_{\text{после}}b_j^2 .
\end{keyeq}
Он опирается на два положения: энергия дислокации пропорциональна $b^2$,
а реакция должна уменьшать энергию. Распад $n$-кратной дислокации~—
частный случай: $n^2b^2>nb^2$.

\begin{vnimanie}
Критерий Франка учитывает только упругие энергии дислокаций. Он не
учитывает ни энергию дефекта упаковки, который может возникнуть при
реакции, ни энергию взаимодействия продуктов реакции. Если суммы
квадратов равны, критерий не даёт ответа.
\end{vnimanie}

\subsubsection{Плотнейшие упаковки и дефекты упаковки}

Чтобы понять частичные дислокации в ГЦК и ГПУ металлах, нужно ясно
представлять себе плотнейшие упаковки (п.~1.4.1). Плотноупакованный
слой шаров~— гексагональная сетка (рис.~\ref{fig:abc-layer}). Обозначим
положения шаров этого слоя буквой $A$. Следующий слой ложится в лунки,
причём лунок двух сортов: $B$ (центры треугольников, направленных
вершиной вверх) и $C$ (вершиной вниз). Занять можно только один сорт
лунок.

\begin{figure}[ht]
\centering
\begin{tikzpicture}[scale=1.12]
  \def\s{1.0}
  \foreach \j in {0,...,3}{\foreach \i in {0,...,5}{
     \pgfmathsetmacro{\xx}{\i*\s+0.5*\s*mod(\j,2)}
     \pgfmathsetmacro{\yy}{\j*0.866*\s}
     \node[circle, draw=atomLine, fill=atomA!55, minimum size=10.2mm, inner sep=0] at (\xx,\yy) {};
     \node[font=\footnotesize] at (\xx,\yy) {$A$};}}
  % положения B и C
  \foreach \j in {0,...,2}{\foreach \i in {0,...,4}{
     \pgfmathsetmacro{\xb}{\i*\s+0.5*\s+0.5*\s*mod(\j,2)}
     \pgfmathsetmacro{\yb}{\j*0.866*\s+0.2887*\s}
     \pgfmathsetmacro{\xc}{\i*\s+0.5*\s*mod(\j,2)}
     \pgfmathsetmacro{\yc}{\j*0.866*\s+0.5774*\s}
     \node[font=\footnotesize\bfseries, text=fmBlue] at (\xb,\yb) {$B$};
     \node[font=\footnotesize\bfseries, text=fmRed] at (\xc+1,\yc) {$C$};}}
  % векторы
  \coordinate (B1) at (1.5,0.2887); \coordinate (B2) at (2.5,0.2887); \coordinate (Cm) at (2,0.5774);
  \draw[-{Stealth[length=2.6mm]}, line width=1.4pt, black] (B1) -- (B2) node[midway, below=1pt, font=\footnotesize]{$\vb b_1$};
  \draw[-{Stealth[length=2.6mm]}, line width=1.4pt, fmGreen] (B1) -- (Cm) node[midway, above left=-1pt, font=\footnotesize]{$\vb b_2$};
  \draw[-{Stealth[length=2.6mm]}, line width=1.4pt, fmGreen] (Cm) -- (B2) node[midway, above right=-1pt, font=\footnotesize]{$\vb b_3$};
\end{tikzpicture}
\caption{Плотноупакованный слой шаров $A$ и два сорта лунок над ним: $B$ и
$C$. Сдвиг шара следующего слоя из лунки $B$ в соседнюю лунку $B$ прямо
(вектор $\vb b_1$) требует перекатиться через шар $A$; гораздо легче
пройти зигзагом через лунку $C$ (векторы $\vb b_2$ и $\vb b_3$)}
\label{fig:abc-layer}
\end{figure}

\begin{itemize}
\item Чередование $ABABAB\ldots$~— \term{гексагональная плотнейшая упаковка}
(ГПУ); положения $C$ не заняты.
\item Чередование $ABCABC\ldots$~— \term{кубическая плотнейшая упаковка}
(ГЦК); плотноупакованные слои~— плоскости $\plfam{111}$.
\end{itemize}

\begin{opredelenie}
\term{Дефект упаковки}~— прослойка с нарушенным (по сравнению с
идеальной структурой) чередованием плотноупакованных слоёв.
\end{opredelenie}

Дефект упаковки можно создать сдвигом в плоскости плотнейшей упаковки,
удалением части плотноупакованного слоя или внедрением лишнего слоя
(рис.~\ref{fig:stacking}). Удобно следить за «тройками» соседних слоёв:
тройка вида $ABC$ (все три буквы разные)~— «кубическая», тройка вида
$ABA$~— «гексагональная».

\begin{figure}[ht]
\centering
\begin{tikzpicture}[scale=0.5]
  % #1 — сдвиг панели, #2 — последовательность (список смещений 0/1/2), #3 — подпись
  \newcommand{\stackpanel}[4]{%
    \begin{scope}[shift={(#1,0)}]
    \foreach \L [count=\k from 0] in {#2}{
      \pgfmathsetmacro{\off}{\L/3}
      \foreach \i in {0,...,3}{\node[atom=3.2mm, fill=#4] at (\i+\off,\k*0.82) {};}
      \pgfmathsetmacro{\lab}{\L}
      \node[font=\scriptsize] at (4.6,\k*0.82) {\ifcase\L A\or B\or C\fi};}
    \draw[fmRed, line width=0.9pt] plot coordinates {#3};
    \end{scope}}
  \stackpanel{0}{0,1,0,1,0,1,0,1}{(1,0) (1.333,0.82) (1,1.64) (1.333,2.46) (1,3.28) (1.333,4.1) (1,4.92) (1.333,5.74)}{atomA}
  \stackpanel{6.5}{0,1,2,0,1,2,0,1}{(1,0) (1.333,0.82) (1.667,1.64) (2,2.46) (2.333,3.28) (2.667,4.1) (3,4.92) (3.333,5.74)}{atomA}
  \stackpanel{13}{0,1,2,0,2,0,1,2}{(1,0) (1.333,0.82) (1.667,1.64) (2,2.46) (1.667,3.28) (2,4.1) (2.333,4.92) (2.667,5.74)}{atomA}
  \stackpanel{19.5}{0,1,2,0,2,1,2,0}{(1,0) (1.333,0.82) (1.667,1.64) (2,2.46) (1.667,3.28) (1.333,4.1) (1.667,4.92) (2,5.74)}{atomA}
  \fill[sf, opacity=0.5] (12.4,2.1) rectangle (17.8,3.6);
  \fill[sf, opacity=0.5] (18.9,2.1) rectangle (24.3,4.4);
  \node[lbl] at (2,-1.3) {а}; \node[lbl] at (8.5,-1.3) {б};
  \node[lbl] at (15,-1.3) {в}; \node[lbl] at (21.5,-1.3) {г};
\end{tikzpicture}
\caption{Чередование плотноупакованных слоёв (вид вдоль плотноупакованного
направления; каждый ряд кружков~— один слой, его сдвиг соответствует
положению $A$, $B$ или $C$): \textit{а}~— ГПУ, $ABAB\ldots$ (зигзаг);
\textit{б}~— ГЦК, $ABCABC\ldots$ (прямая «лестница»);
\textit{в}~— дефект упаковки вычитания в ГЦК: удалён слой $B$,
$\ldots ABCA\,CABC\ldots$; \textit{г}~— дефект упаковки внедрения:
вставлен лишний слой $C$, $\ldots ABCA\,C\,BCA\ldots$ Цветом выделены
области с «гексагональным» чередованием}
\label{fig:stacking}
\end{figure}

\paragraph{Дефекты упаковки в ГЦК-структуре}
\begin{itemize}
\item \emph{Сдвиг.} Если слой $B$ и все вышележащие сдвинуть так, чтобы
атомы попали в соседние лунки ($A\to B$, $B\to C$, $C\to A$):
\[
ABC\,|\,ABCABC\ldots\;\longrightarrow\;ABC\,|\,BCABCA\ldots
\]
Около плоскости сдвига появилось чередование $CBC$~— тонкая прослойка
ГПУ-структуры.
\item \emph{Вычитание.} Изымем слой $B$ и сблизим половинки кристалла:
$\ldots ABCA\,CABC\ldots$ Снова прослойка $CAC$ с гексагональным
чередованием~— \term{дефект упаковки вычитания} (внутренний).
\item \emph{Внедрение.} Вставим лишний слой $C$ между $A$ и $B$:
$\ldots ABCA\,C\,BCAB\ldots$~— \term{дефект упаковки внедрения}
(внешний); в нём две гексагональные тройки подряд.
\end{itemize}

\paragraph{Дефекты упаковки в ГПУ-структуре}
\begin{itemize}
\item \emph{Сдвиг} слоя $A$ и всех вышележащих в соседние лунки:
$ABAB\,|\,ABAB\ldots\;\to\;ABAB\,|\,CACA\ldots$ Около плоскости сдвига~—
тройки $ABC$ и $BCA$, т.\,е. прослойка ГЦК-структуры; выше и ниже~— снова
ГПУ ($ABAB$ и $CACA$).
\item \emph{Вычитание.} Просто изъять слой и сблизить половинки нельзя:
встретятся одинаковые слои ($AA$ или $BB$), и упаковка не будет
плотнейшей. Нужен ещё сдвиг одной половинки, например
$ABABA\,CACA\ldots$~— прослойка ГЦК.
\item \emph{Внедрение} слоя $C$: $ABAB\,C\,ABAB$~— тоже прослойка ГЦК.
\end{itemize}

\begin{glavnoe}
Дефект упаковки в ГЦК-кристалле~— тонкая прослойка ГПУ-структуры, а в
ГПУ-кристалле~— прослойка ГЦК-структуры.
\end{glavnoe}

\paragraph{Энергия дефекта упаковки}
Дефект упаковки~— поверхностный (двумерный) дефект. Число ближайших
соседей и расстояния до них в нём не меняются~— меняется лишь
расположение следующих, более далёких слоёв. Зонная теория показывает,
что из-за этого возрастает энергия электронного газа. Избыточная
свободная энергия единицы площади дефекта~— \term{энергия дефекта упаковки}
$\gamma$ (размерность Дж/м$^2$ = Н/м, как у поверхностного натяжения).
Дефект упаковки стремится уменьшить свою площадь.

\begin{table}[ht]
\centering
\caption{Энергия дефекта упаковки некоторых ГЦК-металлов и сплавов
(справочные данные), мДж/м$^2$}
\label{tab:sfe}
\small
\begin{tabular}{@{}lccccccc@{}}
\toprule
 & Al & Ni & Cu & Au & Ag & сталь 18-8 & латунь Cu--30\,\%Zn \\
\midrule
$\gamma$ & 166 & 125 & 45 & 32 & 16 & $\sim$20 & $\sim$7 \\
\bottomrule
\end{tabular}
\end{table}

\subsubsection{Полные дислокации в ГЦК, ОЦК и ГПУ}

\paragraph{ГПУ}
Плоскость плотнейшей упаковки~— базис $(0001)$, направление плотнейшей
упаковки~— $\nfam{11\br20}$. Наименьший вектор трансляции
$\vb b_1=\frac13\nfam{11\br20}$, его длина равна $a$. Другие полные
векторы: $\vb b_2=\nfam{0001}$ (длина $c$) и $\vb b_3=\frac13\nfam{11\br23}$
(длина $\sqrt{a^2+c^2}$); для идеальной упаковки $c=\sqrt{8/3}\,a\approx
1{,}633\,a$. Скольжение наблюдают по плоскостям базиса $\plfam{0001}$,
призмы $\plfam{10\br10}$, пирамид первого рода $\plfam{10\br11}$ и второго
рода $\plfam{11\br22}$.

Единичная краевая дислокация с вектором $\frac13\nfam{11\br20}$ в базисной
плоскости~— это край \emph{двух} полуплоскостей (рис.~\ref{fig:hcp-edge}):
каждая полуплоскость содержит атомы в положениях $A$ и $B$, и только
вместе они дают сдвиг на полный период.

\paragraph{ГЦК}
Полные дислокации: $\frac a2\nfam{110}$, $a\nfam{100}$, $a\nfam{110}$;
наименьший и самый распространённый вектор~— $\frac a2\nfam{110}$ (длина
$a/\sqrt2$). Линии дислокаций лежат в плоскостях скольжения $\plfam{111}$;
у краевой дислокации $\frac a2\nfam{110}$, как и в ГПУ, две
экстраполуплоскости.

\paragraph{ОЦК}
Полные дислокации: $\frac a2\nfam{111}$, $a\nfam{100}$, $a\nfam{110}$.
Наименьший вектор $\frac a2\nfam{111}$ (длина $a\sqrt3/2$), и дислокации с
ним встречаются чаще всего. Представить геометрию дислокаций в ОЦК
труднее: здесь нет простого чередования плотноупакованных слоёв.
Например, слои $\plfam{112}$ чередуются по шестислойной схеме
$ABCDEFABCDEF\ldots$

\begin{figure}[ht]
\centering
\begin{tikzpicture}[scale=0.55]
  \foreach \i in {0,...,9}{\foreach \j in {-3,...,-1}{\node[atom=3.4mm] at (\i-4.5,\j*0.9) {};}}
  \foreach \i in {-6,...,5}{\foreach \j in {0,...,2}{
    \pgfmathsetmacro{\xs}{(\i+0.5)*9/11}
    \pgfmathtruncatemacro{\hl}{(\i==-1 || \i==0) ? 1 : 0}
    \ifnum\hl=1 \node[atomB=3.4mm] at (\xs,\j*0.9) {};\else\node[atom=3.4mm] at (\xs,\j*0.9) {};\fi}}
  \draw[fmGray, dashed] (-5.8,-0.45) -- (5.8,-0.45);
  \dislT[fmRed]{(0,-0.5)}{0}
  \node[slbl, above] at (-0.41,2.3) {I};
  \node[slbl, above] at (0.41,2.3) {II};
\end{tikzpicture}
\caption{Единичная краевая дислокация $\frac13\nfam{11\br20}$ в ГПУ (или
$\frac a2\nfam{110}$ в ГЦК), схема: над плоскостью скольжения две
«лишние» полуплоскости I и II (оранжевые); вместе они дают сдвиг на полный
период решётки}
\label{fig:hcp-edge}
\end{figure}

\subsubsection{Частичные дислокации Шокли. Расщеплённые дислокации}

Вернёмся к рис.~\ref{fig:abc-layer}. Полной дислокации отвечает сдвиг
шара второго слоя из лунки $B$ в соседнюю лунку $B$ на вектор $\vb b_1$.
Прямой путь проходит через вершину шара $A$~— это высокий энергетический
барьер. Гораздо легче дойти до той же лунки зигзагом: сначала по
«желобку» в лунку $C$ (вектор $\vb b_2$), затем по другому желобку в
лунку $B$ (вектор $\vb b_3$). Положение в промежуточной лунке $C$
механически устойчиво: шар плотно прилегает к трём шарам нижнего слоя.
Но это не абсолютный минимум: положение $C$ «неправильное» для данной
структуры, т.\,е. это дефект упаковки с избыточной энергией $\gamma$
(рис.~\ref{fig:shockley-energy}).

\begin{figure}[ht]
\centering
\begin{tikzpicture}
\begin{axis}[width=0.66\textwidth, height=4.8cm, xmin=0, xmax=1, ymin=0, ymax=1.15,
  xtick={0,0.5,1}, xticklabels={$B$,$C$,$B$}, ytick=\empty,
  xlabel={путь шара второго слоя}, ylabel={$E$}, legend pos=north west, legend cell align=left]
  \addplot[black, dashed, thick, domain=0:1, samples=100] {sin(deg(pi*x))^2};
  \addlegendentry{прямой путь $\vb b_1$}
  \addplot[fmGreen, very thick, domain=0:1, samples=150] {0.38*sin(deg(2*pi*x))^2+0.12*sin(deg(pi*x))^2};
  \addlegendentry{зигзаг $\vb b_2+\vb b_3$}
  \draw[dimline] (axis cs:0.53,0) -- node[right, font=\footnotesize]{$\propto\gamma$} (axis cs:0.53,0.12);
\end{axis}
\end{tikzpicture}
\caption{Энергия при сдвиге шара второго слоя из лунки $B$ в соседнюю
лунку $B$ (схема): зигзагообразный путь через лунку $C$ требует
преодоления двух малых барьеров вместо одного большого; промежуточное
положение $C$~— относительный минимум (дефект упаковки)}
\label{fig:shockley-energy}
\end{figure}

\begin{opredelenie}
\term{Частичная дислокация Шокли}~— граница внутри плоскости
плотнейшей упаковки между областью, где прошёл сдвиг на «половинный»
вектор $\vb b_2$ (и образовался дефект упаковки), и областью, где сдвига
не было. Её вектор Бюргерса лежит в плоскости дефекта упаковки:
в ГПУ $\frac13\nfam{1\br100}$ (длина $a/\sqrt3$), в ГЦК
$\frac a6\nfam{112}$ (длина $a/\sqrt6$).
\end{opredelenie}

Вектор Шокли соединяет вершину равностороннего треугольника со стороной
$a$ (период в плотноупакованном слое) с его центром, отсюда длина
$a/\sqrt3$ (в ГЦК период слоя $a/\sqrt2$, и длина $a/\sqrt6$).

Если дефект упаковки заканчивается внутри кристалла, его ограничивают
частичные дислокации: иначе атомы на краю дефекта оказались бы на
недопустимо малом расстоянии от атомов в исходных положениях.

\begin{opredelenie}
\term{Растянутая} (\term{расщеплённая}) \term{дислокация}~— две
частичные дислокации Шокли, связанные полосой дефекта упаковки
(рис.~\ref{fig:extended}). Она возникает при расщеплении полной
дислокации:
\[
\text{ГПУ: }\tfrac13[\br12\br10]=\tfrac13[01\br10]+\tfrac13[\br1100];\qquad
\text{ГЦК: }\tfrac a2[01\br1]=\tfrac a6[\br12\br1]+\tfrac a6[11\br2].
\]
\end{opredelenie}

Обе реакции выгодны по критерию Франка: в ГПУ $a^2>a^2/3+a^2/3$, в ГЦК
$a^2/2>a^2/6+a^2/6$.

\begin{figure}[ht]
\centering
\begin{tikzpicture}[scale=0.9]
  \fill[sf] (0,0) rectangle (6,1.4);
  \draw[disl] (0,0) -- (6,0);
  \draw[disl] (0,1.4) -- (6,1.4);
  \node[slbl, right] at (6,1.4) {головная частичная ($\vb b_2$)};
  \node[slbl, right] at (6,0) {замыкающая частичная ($\vb b_3$)};
  \node[slbl] at (3,0.7) {дефект упаковки};
  \draw[dimline] (-0.4,0) -- node[left, slbl]{$d_0$} (-0.4,1.4);
  \draw[arr, fmGray] (3,1.6) -- (3,2.3) node[above, slbl]{направление скольжения};
  \begin{scope}[shift={(-3.5,0.2)}]
    \draw[-{Stealth[length=2mm]}, line width=1.1pt] (0,0) -- (0,1) node[midway, left, slbl]{$\vb b$};
    \draw[-{Stealth[length=2mm]}, line width=1.1pt, fmGreen] (0,0) -- (0.29,0.5) node[midway, right, slbl]{$\vb b_3$};
    \draw[-{Stealth[length=2mm]}, line width=1.1pt, fmGreen] (0.29,0.5) -- (0,1) node[midway, right, slbl]{$\vb b_2$};
  \end{scope}
\end{tikzpicture}
\caption{Расщеплённая дислокация (вид на плоскость скольжения): две
частичные дислокации Шокли, между которыми полоса дефекта упаковки
шириной $d_0$. Слева~— сложение векторов Бюргерса $\vb b=\vb b_2+\vb b_3$}
\label{fig:extended}
\end{figure}

\paragraph{Ширина расщепления}
Частичные дислокации одного знака отталкиваются (сила $\propto1/d$), а
дефект упаковки стягивает их, как поверхностное натяжение (сила $\gamma$
на единицу длины). В равновесии
\begin{keyeq}\label{eq:sf-width}
d_0\approx\frac{Gb_p^2}{8\pi\gamma}\cdot\frac{2-\nu}{1-\nu}\left(1-\frac{2\nu\cos2\beta}{2-\nu}\right),
\end{keyeq}
где $b_p$~— длина вектора Шокли, $\beta$~— угол между линией и вектором
Бюргерса полной дислокации. Главное здесь~— зависимость
$d_0\propto Gb^2/\gamma$.

\begin{primer}[Ширина расщепления в меди и алюминии]
Медь: $G=48$~ГПа, $b_p=0{,}148$~нм, $\gamma=45$~мДж/м$^2$, $\nu=0{,}34$.
Для краевой дислокации ($\beta=90^\circ$) $d_0\approx3{,}3$~нм ($\approx13b$),
для винтовой ($\beta=0$) $d_0\approx1{,}4$~нм. Алюминий ($\gamma=166$~мДж/м$^2$):
$d_0\approx0{,}6$~нм~— расщепление всего на пару межатомных расстояний.
\end{primer}

\begin{glavnoe}
Энергия дефекта упаковки~— важнейшая характеристика ГЦК-металла. Чем
меньше $\gamma$, тем шире расщеплены дислокации. Расщеплённой винтовой
дислокации трудно поперечно скользить: частичные дислокации сначала
должны слиться (сжать дефект упаковки), а это стоит энергии. Поэтому в
металлах с низкой $\gamma$ (медь, латунь, аустенитная сталь) скольжение
более «плоское», а упрочнение при деформации сильнее, чем в металлах с
высокой $\gamma$ (алюминий, никель).
\end{glavnoe}

\paragraph{Движение расщеплённой дислокации}
Растянутая дислокация скользит в плоскости дефекта упаковки как единое
целое: головная частичная дислокация создаёт дефект упаковки, замыкающая
его устраняет. Результат тот же, что и при пробеге полной дислокации,
поскольку $\vb b=\vb b_2+\vb b_3$. В ГЦК растянутая дислокация скользит
по направлению плотнейшей упаковки $\nfam{110}$ точно так же, как в ГПУ.
В действительности частичные дислокации не чисто краевые или винтовые:
векторы $\vb b_2$ и $\vb b_3$ составляют с линией углы $\pm30^\circ$
относительно вектора $\vb b$.

В ОЦК-структуре возможно, например, расщепление
$\frac a2[1\br11]=\frac a6[1\br11]+\frac a3[1\br11]$; оно выгодно по
критерию Франка ($3a^2/4>a^2/12+a^2/3$), но реализуется, только если
энергия образующегося дефекта упаковки мала. В ОЦК-металлах энергия
дефектов упаковки велика, и широкие расщеплённые дислокации в них не
наблюдаются.

\subsubsection{Частичные дислокации Франка}

Дефект упаковки можно получить и иначе. Удалим из ГЦК-кристалла
\emph{часть} плотноупакованного слоя (например, диск вакансий,
п.~\ref{ssec:disl-origin}) и сблизим соседние слои, ликвидировав щель.
Получится дефект упаковки вычитания, а его граница внутри кристалла~—
\term{частичная дислокация Франка} (рис.~\ref{fig:frank-partial}).

\begin{figure}[ht]
\centering
\begin{tikzpicture}[scale=0.42]
  % слева — идеальная упаковка ABCABC, справа — слой B (k=4) удалён
  \foreach \L [count=\k from 0] in {0,1,2,0,1,2,0,1}{
    \pgfmathsetmacro{\off}{\L/3}
    \foreach \i in {0,...,6}{
      \pgfmathtruncatemacro{\cut}{(\k==4 && \i>3) ? 1 : 0}
      \pgfmathsetmacro{\dy}{(\k>4) ? -0.82/(1+exp(min(9,-(\i-4.0)*3))) : 0}
      \ifnum\cut=0
        \ifnum\k=4 \node[atomB=3.2mm] at (\i+\off,\k*0.82+\dy) {};
        \else \node[atom=3.2mm] at (\i+\off,\k*0.82+\dy) {};\fi
      \fi}}
  \foreach \L [count=\k from 0] in {A,B,C,A,B,C,A,B}{\node[font=\scriptsize] at (-1.2,\k*0.82) {\L};}
  \foreach \L [count=\k from 0] in {A,B,C,A,{},C,A,B}{\node[font=\scriptsize] at (8.4,\k*0.82-0.0) {\L};}
  \fill[sf, opacity=0.45] (4.1,2.3) rectangle (7.9,4.0);
  \dislT[fmRed]{(4.0,3.28)}{90}
  \node[slbl, right] at (8.9,3.2) {дефект упаковки};
\end{tikzpicture}
\caption{Частичная дислокация Франка в ГЦК-решётке (схема): справа удалена
часть слоя $B$ (оранжевый), и верхние слои опустились, сомкнув щель;
справа чередование $\ldots ABCA\,CAB\ldots$ с дефектом упаковки. Граница
удалённой части слоя~— дислокация Франка с вектором $\frac a3\nfam{111}$,
перпендикулярным плоскости дефекта}
\label{fig:frank-partial}
\end{figure}

\begin{opredelenie}
\term{Частичная дислокация Франка}~— граница дефекта упаковки, созданного
удалением или внедрением части плотноупакованного слоя. В ГЦК её вектор
Бюргерса $\frac a3\nfam{111}$ (длина $a/\sqrt3$) \emph{перпендикулярен}
плоскости дефекта упаковки; дислокация Франка~— краевая.
\end{opredelenie}

Вставив между слоями $\plfam{111}$ неполный атомный слой, получим дефект
упаковки внедрения; его граница~— тоже дислокация Франка с тем же
$\frac a3\nfam{111}$ (атомы соседних слоёв раздвигаются по нормали).

\begin{glavnoe}
Вектор Бюргерса дислокации Франка не лежит в плоскости дефекта упаковки,
поэтому она \vazhno{не может скользить}. Перемещаться она может только
переползанием~— достраивая неполный слой или отбирая атомы с его кромки.
Такие дислокации называют \term{сидячими} (полузакреплёнными), в отличие
от \term{скользящих} (дислокации Шокли и все полные дислокации).
\end{glavnoe}

Плоская петля, внутри которой заключён дефект упаковки, называется
\term{сидячей петлёй Франка}; её вектор Бюргерса перпендикулярен плоскости
петли. Такие петли образуются при захлопывании вакансионных дисков и при
скоплении межузельных атомов под облучением.

В ГПУ-структуре дислокации Франка ограничивают дефекты, полученные
внедрением или удалением неполного слоя $\plfam{0001}$. Например, если в
чередование $ABAB$ между слоями $A$ и $B$ внедрить диск атомов в
положениях $C$, его граница~— дислокация Франка с вектором
$\frac12\nfam{0001}$ (длина $c/2$), перпендикулярным плоскости базиса. В
ОЦК-структуре сидячая чисто краевая частичная дислокация ограничивает
дефект упаковки из двух неполных слоёв, внедрённых в семейство
$\plfam{112}$; её вектор $\frac a3\nfam{112}$ перпендикулярен $\plfam{112}$.

\begin{itogi}
\begin{itemize}
\item Полная дислокация: $\vb b$~— вектор трансляции; частичная: нет, и
за ней остаётся дефект упаковки. Выгодны дислокации с наименьшим $b$.
\item Критерий Франка: реакция идёт, если $\sum b^2$ уменьшается.
\item ГЦК~— $ABCABC$, ГПУ~— $ABAB$; дефект упаковки в ГЦК~— прослойка
ГПУ и наоборот; энергия дефекта упаковки $\gamma$.
\item Шокли: $\frac a6\nfam{112}$ в ГЦК, лежит в плоскости дефекта,
скользящая. Полная дислокация расщепляется на две частичные Шокли; ширина
$d_0\propto Gb^2/\gamma$.
\item Франк: $\frac a3\nfam{111}$, перпендикулярна дефекту, сидячая
(движется только переползанием).
\end{itemize}
\end{itogi}

\begin{proverka}
\item Почему в кристалле преобладают дислокации с наименьшим вектором
Бюргерса?
\item Проверьте по критерию Франка реакцию $\frac a2[110]=\frac a6[211]+\frac a6[12\br1]$.
\item Чем дефект упаковки вычитания отличается от дефекта внедрения?
\item Почему дислокация Франка не может скользить?
\item В каком металле поперечное скольжение затруднено сильнее~— в
алюминии или в меди? Почему?
\end{proverka}
